\documentclass{article}
\usepackage{amsmath,amssymb}
\usepackage[margin=1in]{geometry}

\title{Stat 110 Strategic Practice 5 \& Homework 5: Solutions}
\author{Prof.\ Joe Blitzstein (Department of Statistics, Harvard University)}
\date{Fall 2011}

\begin{document}
\maketitle

\section*{Strategic Practice 5 Solutions}

\subsection*{1. Poisson Distribution and Poisson Paradigm}

\textbf{Problem 1.} The region covers $5$ square inches and is observed for $t$ minutes, and raindrops arrive at an average rate of $20$ per square inch per minute. The expected number of drops hitting the region in $t$ minutes is therefore $\lambda = 20 \times 5 \times t = 100t$. A Poisson distribution is appropriate here because: the drops arrive independently of each other, their arrival times are uniformly spread through time and space, and for small time intervals or small areas the probability of any drop landing there is tiny. These are precisely the conditions under which a Poisson model is warranted.

With $\lambda = 100t$, the probability of zero drops in $3$ seconds (which is $3/60 = 1/20$ of a minute, so $t = 1/20$) is
$$
P(X = 0) = e^{-100/20} \frac{(100/20)^0}{0!} = e^{-5}.
$$

\bigskip
\textbf{Problem 2.} The seating arrangements for the two courses are independent uniform random permutations of $100$ seats. A student "matches" if they end up in the same seat for both courses. Let $N$ be the total number of students who match; this has the same structure as the de Montmort hat-matching problem.

\textbf{(a)} Let $E_j$ be the event that student $j$ occupies the same seat in both courses. By inclusion-exclusion and the symmetry among students,
$$
P\!\left(\bigcup_{j=1}^{100} E_j\right) = \sum_{j=1}^{100}(-1)^{j-1}\binom{100}{j}P\!\left(\bigcap_{k=1}^{j}E_k\right).
$$
The probability that $j$ specific students all match is $(100-j)!/100!$ (the remaining $100-j$ students can sit anywhere). Substituting,
$$
P\!\left(\bigcup_{j=1}^{100}E_j\right) = \sum_{j=1}^{100}\frac{(-1)^{j-1}}{j!},
$$
and therefore
$$
P(N=0) = 1 - \sum_{j=1}^{100}\frac{(-1)^{j-1}}{j!} = \sum_{j=0}^{100}\frac{(-1)^j}{j!}.
$$

\textbf{(b)} For the Poisson approximation: let $I_i$ indicate that student $i$ matches. Each $I_i$ has $P(I_i=1)=1/100$, so $E(N)=100\cdot(1/100)=1$. The indicators are weakly dependent (pairwise, $P(I_i=1, I_j=1)=(1/100)(1/99)\approx(1/100)^2$), so by the Poisson paradigm $N$ is approximately $\text{Pois}(1)$. Thus
$$
P(N=0) \approx e^{-1},
$$
which is consistent with part (a): the exact answer is a partial sum of the Taylor series for $e^x$ evaluated at $x=-1$, which converges rapidly to $e^{-1}$.

\textbf{(c)} Using the same Poisson approximation $N\sim\text{Pois}(1)$,
$$
P(N \ge 2) = 1 - P(N=0) - P(N=1) \approx 1 - e^{-1} - e^{-1} = 1 - 2e^{-1}.
$$

\bigskip
\textbf{Problem 3.} We have $X\sim\text{Pois}(\lambda)$ with $\lambda$ unknown and $\theta = e^{-3\lambda}$.

\textbf{(a)} Using the Poisson PMF and the Taylor series for $e^z$,
$$
E(e^{-3X}) = e^{-\lambda}\sum_{k=0}^{\infty}\frac{e^{-3k}\lambda^k}{k!} = e^{-\lambda}\sum_{k=0}^{\infty}\frac{(e^{-3}\lambda)^k}{k!} = e^{-\lambda}\cdot e^{e^{-3}\lambda} = e^{\lambda(e^{-3}-1)}.
$$
Since $\theta = e^{-3\lambda}$, we have $E(e^{-3X}) = e^{\lambda(e^{-3}-1)} \ne e^{-3\lambda}$ in general, so $T=e^{-3X}$ is biased for $\theta$.

\textbf{(b)} For $g(X) = (-2)^X$,
$$
E((-2)^X) = e^{-\lambda}\sum_{k=0}^{\infty}\frac{(-2)^k \lambda^k}{k!} = e^{-\lambda} e^{-2\lambda} = e^{-3\lambda} = \theta.
$$
So $g(X)=(-2)^X$ is unbiased for $\theta$.

\textbf{(c)} The estimator $g(X)=(-2)^X$ alternates in sign depending on whether $X$ is even or odd. Since $\theta = e^{-3\lambda} > 0$ always, a negative estimate can never be correct and is always farther from $\theta$ than zero would be. A simple improvement is to define
$$
h(X) = \begin{cases} g(X) & \text{if } g(X) \ge 0, \\ 0 & \text{if } g(X) < 0. \end{cases}
$$
Then $|h(X) - \theta| \le |g(X) - \theta|$ always (since clamping a negative value to $0$ moves it closer to the positive target $\theta$), with strict improvement whenever $g(X) < 0$.

Even better: since $\theta = e^{-3\lambda} \in (0,1)$ for $\lambda > 0$, clipping $g(X)$ to lie in $[0,1]$ makes even more sense, as any value outside $[0,1]$ is automatically a worse estimate than its nearest boundary point.

Note: the fact that $g(X)$ is the unique unbiased estimator of $\theta$ doesn't make it a sensible choice in practice. Unbiasedness is just one criterion, and here it forces us to use a clearly irrational estimator.

\subsection*{2. Seeking Sublime Symmetry}

\textbf{Problem 1.} We find the CDF of $SZ$ by conditioning on the value of $S$. For any real $x$,
\begin{align*}
P(SZ \le x) &= P(SZ \le x \mid S=1)\cdot\tfrac12 + P(SZ \le x \mid S=-1)\cdot\tfrac12 \\
&= P(Z \le x)\cdot\tfrac12 + P(-Z \le x)\cdot\tfrac12 \\
&= \Phi(x)\cdot\tfrac12 + P(Z \ge -x)\cdot\tfrac12 \\
&= \Phi(x)\cdot\tfrac12 + \Phi(x)\cdot\tfrac12 = \Phi(x),
\end{align*}
where the last step uses the symmetry $P(Z \ge -x) = P(Z \le x) = \Phi(x)$ for $Z\sim N(0,1)$. Since the CDF of $SZ$ equals $\Phi$, we conclude $SZ\sim N(0,1)$.

\bigskip
\textbf{Problem 2.} When $X$ and $Y$ are i.i.d., swapping their labels changes nothing about the probability model. So replacing $(X,Y)$ by $(Y,X)$ in any probability statement leaves it unchanged, which gives $P(X<Y)=P(Y<X)$.

However, this does not guarantee $P(X<Y)=1/2$. In the discrete case, $P(X=Y)>0$, so $P(X<Y)+P(Y<X)+P(X=Y)=1$ gives $P(X<Y)=P(Y<X)<1/2$. In the continuous case, $P(X=Y)=0$, so the two probabilities do sum to $1$, giving $P(X<Y)=1/2$.

If $X$ and $Y$ merely share a distribution but are not independent, the symmetry argument breaks down. As a simple example, let $X$ always equal $Y$ (total dependence): then $P(X<Y)=0$ and $P(Y<X)=0$, so the probabilities are still equal here, but other dependence structures can make them unequal. For instance, if $X$ and $Y$ have the same marginal but $Y=X+c$ for some positive constant $c$ almost surely, then $P(X<Y)=1$ but $P(Y<X)=0$.

\bigskip
\textbf{Problem 3.} If $X\sim\text{Bin}(n,p)$, think of $X$ as the number of successes in $n$ independent trials with success probability $p$. Then $n-X$ is the number of failures. Relabeling failure as success (and vice versa) gives an equivalent experiment where each trial succeeds (formerly failed) with probability $1-p$. So $n-X\sim\text{Bin}(n,1-p)$.

\bigskip
\textbf{Problem 4.} Number the seats so that seat $j$ is the assigned seat for passenger $j$. Consider what seats could possibly be open for the last passenger (passenger $100$). Any intermediate seat $j$ (for $2\le j\le 99$) is always taken by passenger $j$ if available, so by the time passenger $100$ boards, seat $j$ is either taken by passenger $j$ or by a prior passenger who was displaced. The key observation is: when each displaced passenger picks a random seat, the only seats that remain "in contention" throughout the entire boarding process are seat $1$ (the crazy first passenger's assigned seat) and seat $100$ (the last passenger's seat). Seat $j$ for $2\le j\le99$ gets resolved the moment passenger $j$ boards.

By symmetry, every intermediate displaced passenger views seat $1$ and seat $100$ as equally likely candidates when they must choose randomly, since no prior information distinguishes one from the other. Therefore seat $1$ and seat $100$ are equally likely to be the last unresolved seat when passenger $100$ boards, each with probability $1/2$. So
$$
P(\text{last passenger gets seat } 100) = \frac{1}{2}.
$$

\subsection*{3. Continuous Distributions}

\textbf{Problem 1.} Since $Y=e^X$ with $X\sim N(\mu,\sigma^2)$, we have $Y>0$ always. For $y>0$,
$$
P(Y\le y) = P(e^X \le y) = P(X \le \ln y) = P\!\left(\frac{X-\mu}{\sigma}\le\frac{\ln y - \mu}{\sigma}\right) = \Phi\!\left(\frac{\ln y - \mu}{\sigma}\right).
$$
The CDF is $\Phi\!\left(\frac{\ln y-\mu}{\sigma}\right)$ for $y>0$ and $0$ for $y\le 0$. Differentiating with respect to $y$ via the chain rule,
$$
f_Y(y) = \frac{1}{\sigma y\sqrt{2\pi}}\exp\!\left(-\frac{(\ln y - \mu)^2}{2\sigma^2}\right), \quad y>0.
$$

\bigskip
\textbf{Problem 2.} The desired distribution has CDF $F(x)=\int_0^x \lambda e^{-\lambda t}\,dt = 1-e^{-\lambda x}$ for $x>0$. To generate $X$ from $U\sim\text{Unif}(0,1)$, invert the CDF: set $F(X)=U$, giving $1-e^{-\lambda X}=U$, so
$$
X = F^{-1}(U) = -\frac{1}{\lambda}\ln(1-U).
$$
Since $1-U$ is also $\text{Unif}(0,1)$, this simplifies to $X = -\frac{1}{\lambda}\ln(U)$. By the Universality of the Uniform, $X$ has CDF $F$, i.e., $X\sim\text{Exponential}(\lambda)$.

\bigskip
\textbf{Problem 3.} By Universality of the Uniform, for any continuous random variable $W$ with CDF $G$, we have $G(W)\sim\text{Unif}(0,1)$. Taking $W=Z\sim N(0,1)$ with CDF $\Phi$, it follows that $\Phi(Z)\sim\text{Unif}(0,1)$. Since the mean of a $\text{Unif}(0,1)$ is $1/2$,
$$
E(\Phi(Z)) = \frac{1}{2}.
$$

\bigskip
\textbf{Problem 4.} Scale so the stick has length $1$, and let the break point $U\sim\text{Unif}(0,1)$. The longer piece has length $L=\max(U,1-U)$. For $l\in[1/2,1]$, the event $\{L\le l\}$ is equivalent to $\{U\le l\}\cap\{1-U\le l\}$, which simplifies to $\{1-l\le U\le l\}$. So
$$
F_L(l) = P(L\le l) = P(1-l\le U\le l) = 2l-1, \quad l\in[1/2,1].
$$
This is the CDF of a $\text{Unif}(1/2,1)$ distribution, so $L\sim\text{Unif}(1/2,1)$ and
$$
E(L)=\frac{1/2+1}{2}=\frac{3}{4}.
$$

\subsection*{4. LOTUS}

\textbf{Problem 1.} For $X\sim\text{Pois}(\lambda)$, applying LOTUS,
$$
E(X!) = \sum_{k=0}^{\infty} k!\cdot\frac{e^{-\lambda}\lambda^k}{k!} = e^{-\lambda}\sum_{k=0}^{\infty}\lambda^k.
$$
This is a geometric series in $\lambda$: it converges to $e^{-\lambda}\cdot\frac{1}{1-\lambda}$ when $0<\lambda<1$, and diverges for $\lambda\ge1$. So
$$
E(X!) = \frac{e^{-\lambda}}{1-\lambda}, \quad 0<\lambda<1.
$$

\bigskip
\textbf{Problem 2.} By LOTUS with the $N(0,1)$ PDF $\varphi(z)=\frac{1}{\sqrt{2\pi}}e^{-z^2/2}$,
$$
E|Z| = \int_{-\infty}^{\infty}|z|\varphi(z)\,dz = 2\int_0^{\infty}z\varphi(z)\,dz = \frac{2}{\sqrt{2\pi}}\int_0^{\infty}ze^{-z^2/2}\,dz.
$$
Substituting $u=z^2/2$ (so $du=z\,dz$),
$$
E|Z| = \frac{2}{\sqrt{2\pi}}\int_0^{\infty}e^{-u}\,du = \frac{2}{\sqrt{2\pi}} = \sqrt{\frac{2}{\pi}}.
$$

\bigskip
\textbf{Problem 3.} Let $q=1-p$ and $X\sim\text{Geom}(p)$, so $P(X=k)=pq^k$ for $k=0,1,2,\dots$. By LOTUS,
$$
E(e^{tX}) = \sum_{k=0}^{\infty}e^{tk}\cdot pq^k = p\sum_{k=0}^{\infty}(qe^t)^k = \frac{p}{1-qe^t},
$$
valid as long as $qe^t<1$, i.e., $t<-\ln q = \ln(1/q)$. The series diverges for $qe^t\ge1$.

\bigskip
\section*{Homework 5 Solutions}

\subsection*{1. Birthday Indicators}

\textbf{(a)} $I_{12}$ and $I_{34}$ are independent: they depend on entirely different people, so knowing the birthdays of persons $1$ and $2$ tells us nothing about persons $3$ and $4$.

$I_{12}$ and $I_{13}$ are also independent, even though both involve person $1$: learning that persons $1$ and $2$ share a birthday shifts our knowledge about person $2$'s birthday but, since birthdays are uniform over $365$ days, leaves $P(\text{person 1's birthday} = \text{person 3's birthday})$ unchanged at $1/365$.

The $I_{ij}$'s as a whole collection are \emph{not} mutually independent, however. A simple illustration: if $I_{12}=1$ and $I_{13}=1$ (persons $1,2,3$ all share a birthday), then $I_{23}=1$ is forced. So the joint event $\{I_{12}=1, I_{13}=1, I_{23}=0\}$ has probability zero, yet each indicator individually has positive probability of being $1$, contradicting mutual independence.

\bigskip
\textbf{(b)} Let $X$ be the number of birthday-matching pairs. There are $\binom{n}{2}$ pairs, each matching with probability $1/365$, so $E(X)=\binom{n}{2}/365$. The pairs are pairwise independent (as shown in (a)) and there are a large number of them each with small probability, which are the conditions for the Poisson paradigm. So $X\approx\text{Pois}\!\left(\frac{n(n-1)}{730}\right)$.

For $n=23$, $\lambda=\frac{23\cdot22}{730}\approx0.693$. Then
$$
P(X\ge1)\approx 1-e^{-0.693}\approx 0.500.
$$

\bigskip
\textbf{(c)} Now redefine $I_{ij}=1$ if persons $i$ and $j$ have the same birthday or birthdays one day apart. Two birthdays are "within one day" in $3$ out of $365$ ways (same day, one day before, one day after, accounting for the cyclical calendar), so $P(I_{ij}=1)=3/365$. By the Poisson paradigm, the number of near-matching pairs is approximately $\text{Pois}\!\left(\frac{3n(n-1)}{730}\right)$. Setting $P(X\ge1)\approx1/2$,
$$
1-e^{-3n(n-1)/730}=\frac{1}{2} \implies \frac{3n(n-1)}{730}=\ln2 \implies 3n^2-3n-730\ln2\approx0.
$$
Solving the quadratic (taking the positive root) gives $n\approx13.5$, so $n=14$ people suffice. The exact computation also yields $14$, confirming the accuracy of the approximation.

\subsection*{2. Discretizing an Exponential}

$T$ has PDF $\frac{1}{5}e^{-t/5}$ and CDF $F_T(t)=1-e^{-t/5}$ for $t>0$. The event $\{X=k\}$ corresponds to the news being posted on day $k+1$, i.e., $\frac{k}{365}\le T < \frac{k+1}{365}$. So
$$
P(X=k) = F_T\!\left(\frac{k+1}{365}\right) - F_T\!\left(\frac{k}{365}\right) = e^{-k/1825} - e^{-(k+1)/1825}.
$$
Factoring out $e^{-k/1825}$,
$$
P(X=k) = e^{-k/1825}\!\left(1-e^{-1/1825}\right) = (1-p)^k p, \quad k=0,1,2,\dots,
$$
where $p=1-e^{-1/1825}$. This is exactly the $\text{Geometric}(p)$ distribution (with support on $\{0,1,2,\dots\}$). The small value $p\approx0.0005$ reflects both the long timescale of the exponential (mean $5$ years $= 1825$ days) and the fact that the news is unlikely to appear on any given individual day.

\subsection*{3. Uniform on $(-1,1)$}

\textbf{(a)} The PDF of $U$ is $f(u)=1/2$ for $u\in(-1,1)$. By symmetry, $E(U)=0$. Using LOTUS,
$$
E(U^2)=\frac{1}{2}\int_{-1}^{1}u^2\,du=\frac{1}{2}\cdot\frac{2}{3}=\frac{1}{3}.
$$
So $\text{Var}(U)=E(U^2)-(EU)^2=\frac{1}{3}$. Similarly,
$$
E(U^4)=\frac{1}{2}\int_{-1}^{1}u^4\,du=\frac{1}{2}\cdot\frac{2}{5}=\frac{1}{5}.
$$

\textbf{(b)} Since $0\le U^2\le1$, the CDF of $V=U^2$ is $0$ for $t\le0$ and $1$ for $t\ge1$. For $t\in(0,1)$,
$$
G(t)=P(U^2\le t)=P(-\sqrt{t}\le U\le\sqrt{t})=\frac{2\sqrt{t}}{2}=\sqrt{t}.
$$
Differentiating, the PDF is $g(t)=\frac{1}{2\sqrt{t}}=\frac{1}{2}t^{-1/2}$ for $t\in(0,1)$.

This is not a uniform distribution: the PDF is not constant (it blows up near $t=0$). It is in fact a $\text{Beta}(1/2,1)$ distribution, or equivalently related to an arcsine distribution.

\subsection*{4. Area Under the Quantile Function}

We want to evaluate $\int_0^1 F^{-1}(u)\,du$. Let $U\sim\text{Unif}(0,1)$ and define $X=F^{-1}(U)$. By the Universality of the Uniform, $X$ has CDF $F$, so $E(X)=\mu$. Applying LOTUS to the function $F^{-1}$ evaluated at $U$,
$$
\int_0^1 F^{-1}(u)\,du = E(F^{-1}(U)) = E(X) = \mu.
$$
So the area under the quantile function over $[0,1]$ equals the mean $\mu$ of the distribution.

Quick check: for $F$ the $\text{Unif}(0,1)$ CDF, $F^{-1}(u)=u$, so the area is $\int_0^1 u\,du=1/2$, which equals the mean of $\text{Unif}(0,1)$.

\subsection*{5. Censored Normal}

$X = Z\cdot\mathbf{1}_{Z>0}$ where $Z\sim N(0,1)$. So $X=0$ when $Z\le0$ (probability $1/2$) and $X=Z$ when $Z>0$ (probability $1/2$). Applying LOTUS with the $N(0,1)$ PDF $\varphi(z)=\frac{1}{\sqrt{2\pi}}e^{-z^2/2}$,
$$
E(X) = \int_{-\infty}^{\infty}z\cdot\mathbf{1}_{z>0}\cdot\varphi(z)\,dz = \int_0^{\infty}z\varphi(z)\,dz = \frac{1}{\sqrt{2\pi}}\int_0^{\infty}ze^{-z^2/2}\,dz.
$$
Substituting $u=z^2/2$,
$$
E(X) = \frac{1}{\sqrt{2\pi}}\int_0^{\infty}e^{-u}\,du = \frac{1}{\sqrt{2\pi}}.
$$

For the variance, first compute $E(X^2)$. Since $X^2=Z^2\cdot\mathbf{1}_{Z>0}$,
$$
E(X^2) = \int_0^{\infty}z^2\varphi(z)\,dz = \frac{1}{2}\int_{-\infty}^{\infty}z^2\varphi(z)\,dz = \frac{1}{2}E(Z^2) = \frac{1}{2},
$$
where we used symmetry of $\varphi$ and the fact that $\text{Var}(Z)=E(Z^2)=1$. Therefore,
$$
\text{Var}(X) = E(X^2) - (E(X))^2 = \frac{1}{2} - \frac{1}{2\pi}.
$$

Note that $X$ is a mixed random variable: it places mass $1/2$ at $0$ and is otherwise continuous on $(0,\infty)$. It has neither a pure PMF nor a pure PDF, but LOTUS still applies by working through the distribution of $Z$.

As a sanity check: $\text{Var}(X)<1=\text{Var}(Z)$ makes sense since pinning $X$ to $0$ half the time and restricting it to nonneg values both reduce its spread compared to $Z$.

\end{document}